\begin{definition}{Absolute Value $\lvert \cdot \rvert$}{absolute_value}
    If $x$ is a rational number, the absolute value of $\lvert x \rvert$ if $x$ is defined as follows:
    If $x$ is negative, then $\lvert x \rvert := -x.$
    If $x$ is zero, then $\lvert x \rvert := 0.$
    If $x$ is positive, then $\lvert x \rvert := x.$
\end{definition}

\begin{definition}{Distance $d(\cdot, \cdot)$}{distance}
    Let $x$ and $y$ be rational numbers. 
    The quantity $\lvert x - y \rvert$ is called the distance between $x$ and $y$ and is sometimes denoted as $d(x, y) := \lvert x - y \rvert.$
\end{definition}

\begin{definition}{Exponentiation to an integer}{exponentiation}
    Let $x \in \mathbb{Q}$.
    To raise $x$ to the power 0, we define $x^0 := 1;$ in particular, we define $0^0 := 1.$
    Now, suppose inductively that $x^n$ has been defined for some $n \in \mathbb{N}.$
    Then, we define $x^{n+1} := x^n \times x.$

    If instead we wish to raise $x$ to the power of a negative number, we need that $x$ is a non-zero rational and have a negative integer $-n.$
    We define $x^{-n} := \frac{1}{x^n}.$
    
\end{definition}

\begin{lemma}
    \label{lemma: abs_square}
    We have that $|x|^2 = x^2$.
\end{lemma}
\begin{proof}
    Let $x$ to be a rational number. 
    We consider three cases by the law of trichotomy on the rationals.
    \begin{enumerate}
        \item Say $x$ is negative, so $\lvert x \rvert ^2 = (-x)^2 = x^2.$
        \item Say $x = 0$. Observe $\lvert x \rvert^2 = \lvert 0 \rvert^2 =0^2.$
        \item Say $x$ is positive. Then, $\lvert x \rvert^2 = x^2$.
    \end{enumerate}
    In all of the three cases, we see that $\lvert x \rvert ^2 = x^2.$
\end{proof}

% Problem 1
\begin{problem}{Properties of absolute value and distance}
	Verify the basic properties of absolute value and distance. Let $x, y, z \in \mathbb{Q}.$
\end{problem}
%\begin{multicols}{2}
\begin{proposition}
    Non-degeneracy of absolute value. We have $\lvert x \rvert \geq 0.$ 
    Also, $\lvert x \rvert = 0$ if and only if $x = 0.$
\end{proposition}
\begin{proof}
    First, we show that $\lvert x \rvert \geq 0.$ 
    By the law of trichotomy, we consider three cases. 
    If $x$ is negative, then there exists a positive rational number $x'$ such that $x = -(x').$
    Then, $\lvert x \rvert = \lvert -(x') \rvert = x'$ (by multiplicativity of absolute value).
    Since $x'$ is positive, it follows that $\lvert x \rvert \geq 0.$
    If $x = 0$, then clearly $\lvert x \rvert = \lvert 0 \rvert = 0 \geq 0.$
    Finally, if $x$ is positive, then $\lvert x \rvert = x \geq 0.$

    Now, we show that $\lvert x \rvert = 0$ if and only if $x = 0.$
    Now, suppose that $x \neq 0.$ We now consider two cases.
    If $x$ is negative, then $\lvert x \rvert = -x$ is positive. 
    By the law of trichotomy, $\lvert x \rvert$ cannot be 0.
    If $x$ is instead positive, then $\lvert x = \rvert = x$ is also positive, which by the law of trichotomy, cannot be 0.
    Hence, in either case, $\lvert x \rvert \neq 0.$
    Conversely, suppose that $x = 0.$ Trivially, $\lvert x \rvert = \lvert 0 \rvert = 0.$
\end{proof}

\begin{proposition}
    (Triangle inequality for absolute value) We have $\lvert x + y \rvert \leq \lvert x \rvert + \lvert y \rvert.$
\end{proposition}
\begin{proof}
    We wish to show $\lvert x + y \rvert \leq \lvert x \rvert + \lvert y \rvert.$ 
    By definition of order on the rationals, we need to show that 
    $$\lvert x + y \rvert - (\lvert x \rvert + \lvert y \rvert) \leq  0.$$
    To prove this, we now consider two cases of $x + y = 0$ and $x + y \neq 0.$
    \begin{enumerate}
        \item Say $x + y = 0.$ Then, $\lvert x + y \rvert - (\lvert x \rvert + \lvert y \rvert) = 0 - (\lvert x \rvert + \lvert y \rvert) = - (\lvert x \rvert + \lvert y \rvert) \leq 0.$
        Therefore, $\lvert x + y \rvert - (\lvert x \rvert + \lvert y \rvert) \leq 0.$
        \item Say instead that $x + y \neq 0,$ so $x \neq 0$ and $y \neq 0.$ 
        Then, we will use Lemma~\ref{lemma: abs_square} whenever we have $\lvert x \rvert ^2 = x^2.$ Notice:
        \begin{align*}
            \lvert x + y \rvert - (\lvert x \rvert + \lvert y \rvert) \leq 0  &\iff \big( \lvert x + y \rvert - (\lvert x \rvert + \lvert y \rvert) \big) \big( \lvert x + y\rvert + (\lvert x \rvert + \lvert y \rvert) \big) \leq 0 \\
            &\iff \lvert x + y \rvert ^2 - \big(\lvert x \rvert + \lvert y \rvert \big)^2 \leq 0 \\
            &\iff (x + y)^2 - (\lvert x \rvert ^2 + 2 \lvert x \rvert  \lvert y \rvert + \lvert y \rvert ^2) \leq 0 \\
            &\iff x^2 + 2xy + y^2 - (x^2 + 2 \lvert x \rvert  \lvert y \rvert + y^2) \leq 0 \\
            &\iff 2xy - 2 \lvert x \rvert \lvert y \rvert \leq 0 \\
            &\iff xy - \lvert x \rvert \lvert y \rvert \leq  0
        \end{align*}
    We will now just use cases to show that $xy - \lvert x \rvert \lvert y \rvert \leq 0.$ 
    (We can also be smarter by using work that we will do because Propositions 1.3 and 1.4 give us the tools to skip the casework.)
    \begin{enumerate}
        \item Say $x < 0$ and $ y < 0,$ so there exist positive rational numbers $x'$ and $y'$ such that $x = -(x')$ and $y = -(y').$
        Then $xy - \lvert x \rvert \lvert y \rvert = (-x')(-y') - \lvert -x' \rvert \lvert -y' \rvert = x'y' - x'y' = 0 \leq 0.$
        \item Say $x < 0$ and $y > 0,$ so there exists a positive rational number $x'$ such that $x = (-x').$
        Then, $xy - \lvert x \rvert \lvert y \rvert = (-x')y - \lvert -x' \rvert \lvert y \rvert = -x'y -x'y = -2x'y = -2 (x'y) < 0 \leq 0$
        \item Say $x > 0$ and $y <0$. The logic is the same as for the proof above because of symmetry between the $x$ and $y$. In any case, we will have that $xy - \lvert x \rvert \lvert y \rvert < 0 \leq 0$.
        \item Say $x > 0$ and $y > 0,$ so $xy - \lvert x \rvert \lvert y \rvert = xy - xy = 0 \leq 0.$
    \end{enumerate} 
    
    In all of the four cases, we have that $xy - \lvert x \rvert \lvert y \rvert \leq  0.$ 
    Therefore, we can say that $\lvert x + y \rvert - (\lvert x \rvert + \lvert y \rvert) \leq 0$.
    \end{enumerate}
    Overall, in either case of $x + y = 0$ or $x + y \neq 0,$ we have shown that $\lvert x + y \rvert \leq (\lvert x \rvert + \lvert y \rvert)$
\end{proof}

\begin{proposition}
    We have the inequalities $-y \leq x \leq y$ if and only if $y \geq \lvert x \rvert.$
    In particular, we have $-\lvert x \rvert \leq x \leq \lvert x \rvert$.
\end{proposition}
\begin{proof}
    s
\end{proof}

\begin{proposition}
    (Multiplicativity of absolute value) We have $\lvert xy \rvert = \lvert x \rvert \lvert y \rvert.$ 
    In particular, $\lvert -x \rvert = \lvert x \rvert.$
\end{proposition}
\begin{proof}
    s
\end{proof}

\begin{proposition}
    (Non-degeneracy of distance) We have $d(x, y) \geq 0.$
    Also, $d(x, y) = 0$ if and only if $x = y.$
\end{proposition}
\begin{proof}
    smth
\end{proof}

\begin{proposition}
    (Symmetry of distance) We have $d(x, y) = d(y, x).$
\end{proposition}
\begin{proof}
    sh 
\end{proof}

\begin{proposition}
    (Triangle inequality for distance) We have $d(x, z) \leq d(x, y) + d(y, z).$
\end{proposition}

\begin{proof}
    chu
\end{proof}
%\end{multicols}

\begin{definition}{$\varepsilon$-closeness}{closeness}
    Let $\varepsilon > 0$ be a rational number, and let $x, y$ be rational numbers.
    We say that $y$ is $\varepsilon$-close to $x$ if and only if we have that $d(x, y) \leq \varepsilon.$
\end{definition}
% Problem 2
\begin{problem}{Properties of $\epsilon$-closeness}
Let $x, y, z, w \in \mathbb{Q}$. Then, the following properties hold:
\end{problem}
\begin{multicols}{2}
\begin{proposition}
    The number $x$ is $\varepsilon$-close to $y$ for every $\varepsilon > 0$ if and only if  $x = y$.
\end{proposition} 
\begin{proof}
    d
\end{proof}

\begin{proposition}
    Let $\varepsilon > 0.$ 
    If $x$ is $\varepsilon$-close to $y$, then $y$ is $\varepsilon$-close to $x$.
\end{proposition}
\begin{proof}
    s
\end{proof}

\begin{proposition}
    Let $\varepsilon, \delta > 0.$ If $x$ is $\varepsilon$-close to $y$, and $y$ is $\delta$-close to $z$, then $x$ and $z$ are $(\varepsilon + \delta)$-close. 
\end{proposition}
\begin{proof}
    chuchu
\end{proof}

\begin{proposition}
    Let $\varepsilon > 0.$ If $x$ and $y$ are $\varepsilon$-close, they are also $\varepsilon'$-close for every $\varepsilon' > \varepsilon$.
\end{proposition}
\begin{proof}
    chuchu
\end{proof}

\begin{proposition}
    Let $\varepsilon > 0.$ If $y$ and $z$ are both $\varepsilon$-close to $x$ and $w$ is between $y$ and $z$ (i.e., $y \leq w \leq z$ or $z \leq w \leq y$), then $w$ is also $\varepsilon$-close to $x$.
\end{proposition}
\begin{proof}
    chuchu
\end{proof}

\begin{proposition}
    Let $\varepsilon > 0.$ If $x$ and $y$ are $\varepsilon$-close, and $z$ is non-zero, then $xz$ and $yz$ are $\varepsilon \lvert z \rvert$-close.
\end{proposition}
\begin{proof}
    
\end{proof}

\begin{proposition}
    Let $\varepsilon, \delta > 0.$ Suppose $x$ and $y$ are $\varepsilon$-close, and $z$ and $w$ are $\delta$-close. Then $xz$ and $yw$ are $(\varepsilon \lvert z \rvert + \delta \lvert x \rvert + \varepsilon \delta)$-close.
\end{proposition}
\end{multicols}

% Probelm 3
\begin{problem}{Properties of Exponentiation I}
Let $x, y$ be rational numbers. Let $n, m$ be natural numbers. The following properties hold.
\end{problem}

\begin{multicols}{2}
\begin{proposition} We have $x^n x^m = x^{n + m}$.
\end{proposition}
\begin{proof}
    chuchu
\end{proof}

\begin{proposition} We have $(x^n)^m=x^{nm}.$
\end{proposition}
\begin{proof}
    chuchu
\end{proof}

\begin{proposition} $(xy)^n=x^n y^n.$
\end{proposition}
\begin{proof}
    chuchu
\end{proof}
\begin{proposition}
    Suppose $n \geq 0.$ Then, we have $x^n = 0$ if and only if $x = 0.$
\end{proposition}
\begin{proof}
    suppose
\end{proof}

\begin{proposition}
    If $x \geq y \geq 0$, then $x^n \geq y ^ n \geq 0$. 
\end{proposition}
\begin{proof}
    suppose
\end{proof}

\begin{proposition}
    If $x > y \geq 0$ and $n > 0$, then $x ^n > y^n \geq 0.$
\end{proposition}
\begin{proof}
    suppose
\end{proof}
\begin{proposition}
    We have $\lvert x^n \rvert = \lvert x \rvert ^n$.
\end{proposition}
\begin{proof}
    suppose
\end{proof}
\end{multicols}


% Probelm 5
\begin{problem}{Properties of Exponentiation II}
Let $x, y$ be rational numbers. Let $n, m$ be integers. The following properties hold.
\end{problem}

\begin{multicols}{2}
\begin{proposition} We have $x^n x^m = x^{n + m}$.
\end{proposition}
\begin{proof}
    chuchu
\end{proof}

\begin{proposition} We have $(x^n)^m=x^{nm}.$
\end{proposition}
\begin{proof}
    chuchu
\end{proof}

\begin{proposition} $(xy)^n=x^n y^n.$
\end{proposition}
\begin{proof}
    chuchu
\end{proof}

\begin{proposition}
    If $x \geq y > 0$, then $x^n \geq y^n > 0$ if $n$ is positive OR $0 < x^n \leq y^n$ if $n$ is negative.
\end{proposition}
\begin{proof}
    suppose
\end{proof}

\begin{proposition}
    If $x, y > 0,$ $n \neq 0,$ and $x^n = y^n$, then $x = y$.
\end{proposition}
\begin{proof}
    suppose
\end{proof}

\begin{proposition}
    We have $\lvert x^n \rvert = \lvert x \rvert ^n$.
\end{proposition}
\begin{proof}
    suppose
\end{proof}
\end{multicols}

\begin{problem}{Exponentiation grows faster than linear}
    Prove that $2^N \geq N$ for all positive integers $N$.
\end{problem}
\begin{proof}
    chuchu
\end{proof}